\chapter{Complex Numbers, Waves, and Fourier Analysis}

\section{The imaginary unit}

A complex number is \(z=a+\ii b\), where \(a,b\in\R\) and \(\ii^2=-1\).  Its complex conjugate is \(z^*=a-\ii b\), and
\[
  |z|^2=z^*z=(a-\ii b)(a+\ii b)=a^2+b^2.
\]
Write \(z=re^{\ii\theta}\) with
\[
  r=|z|,\qquad a=r\cos\theta,\qquad b=r\sin\theta.
\]
Euler's formula
\[
  e^{\ii\theta}=\cos\theta+\ii\sin\theta
\]
follows by comparing the power series:
\[
  e^{\ii\theta}
  =
  1+\ii\theta-\frac{\theta^2}{2!}
  -\ii\frac{\theta^3}{3!}
  +\frac{\theta^4}{4!}+\cdots .
\]
The real terms form \(\cos\theta\), the imaginary terms form \(\ii\sin\theta\).

\section{Plane waves}

The real wave \(\cos(kx-\omega t)\) is often written as the real part of
\[
  e^{\ii(kx-\omega t)}.
\]
This is not a trick.  Differentiation of an exponential is algebra:
\[
  \frac{\p}{\p t}e^{\ii(kx-\omega t)}=-\ii\omega e^{\ii(kx-\omega t)},\qquad
  \frac{\p}{\p x}e^{\ii(kx-\omega t)}=\ii k e^{\ii(kx-\omega t)}.
\]
Substituting into the wave equation
\[
  \frac{\p^2 u}{\p t^2}-c^2\frac{\p^2u}{\p x^2}=0
\]
gives
\[
  (-\omega^2+c^2k^2)e^{\ii(kx-\omega t)}=0.
\]
Therefore nonzero waves obey the dispersion relation
\[
  \omega^2=c^2k^2.
\]
Field theory repeatedly turns differential equations into algebraic equations by this substitution.

\section{Fourier series}

Suppose \(f(x)\) is periodic with period \(L\).  The basic waves are
\[
  e^{\ii k_nx},\qquad k_n=\frac{2\pi n}{L},\qquad n\in\Z.
\]
They obey the orthogonality relation
\[
  \int_0^L e^{-\ii k_mx}e^{\ii k_nx}\dd x
  =
  \int_0^L e^{\ii(k_n-k_m)x}\dd x
  =
  L\delta_{mn}.
\]
Here \(\delta_{mn}\) is \(1\) if \(m=n\) and \(0\) otherwise.  If
\[
  f(x)=\sum_{n=-\infty}^{\infty}c_ne^{\ii k_nx},
\]
then multiplying by \(e^{-\ii k_mx}\), integrating, and using orthogonality gives
\[
  c_m=\frac1L\int_0^L f(x)e^{-\ii k_mx}\dd x.
\]

\begin{example}[A square-integrable sawtooth]
Let \(f(x)=x\) on \((-\pi,\pi)\), extended periodically.  Since \(f\) is odd,
\[
  f(x)=\sum_{n=1}^{\infty}b_n\sin nx.
\]
The coefficients are
\[
  b_n=\frac1\pi\int_{-\pi}^{\pi}x\sin nx\,\dd x
  =\frac2\pi\int_0^\pi x\sin nx\,\dd x.
\]
Integrating by parts,
\[
  \int_0^\pi x\sin nx\,\dd x
  =
  \left[-\frac{x\cos nx}{n}\right]_0^\pi
  +\frac1n\int_0^\pi\cos nx\,\dd x
  =
  -\frac{\pi(-1)^n}{n}.
\]
Thus
\[
  x=2\sum_{n=1}^\infty \frac{(-1)^{n+1}}{n}\sin nx.
\]
\end{example}

\section{Fourier transforms}

As \(L\to\infty\), allowed \(k\)-values become dense.  The Fourier transform convention used here is
\[
  \tilde f(k)=\int_{-\infty}^{\infty}f(x)e^{-\ii kx}\dd x,
  \qquad
  f(x)=\int_{-\infty}^{\infty}\frac{\dd k}{2\pi}\tilde f(k)e^{\ii kx}.
\]
In three dimensions,
\[
  \tilde f(\bm k)=\int \dd^3x\,f(\bm x)e^{-\ii\bm k\cdot\bm x},
  \qquad
  f(\bm x)=\int\frac{\dd^3k}{(2\pi)^3}\tilde f(\bm k)e^{\ii\bm k\cdot\bm x}.
\]
The Laplacian becomes multiplication:
\[
  \nabla^2 e^{\ii\bm k\cdot\bm x}=-|\bm k|^2e^{\ii\bm k\cdot\bm x}.
\]

\section{The delta function}

The delta function is defined by what it does under an integral:
\[
  \int_{-\infty}^{\infty}\delta(x-a)f(x)\dd x=f(a).
\]
It is not an ordinary function.  A useful approximation is
\[
  \delta_\epsilon(x)=\frac{1}{\sqrt{\pi}\epsilon}e^{-x^2/\epsilon^2}.
\]
Its area is \(1\), and as \(\epsilon\to0\) it becomes narrow.  For a smooth \(f\),
\[
  \int \delta_\epsilon(x-a)f(x)\dd x
  =
  \int \frac{e^{-u^2/\epsilon^2}}{\sqrt{\pi}\epsilon}
  f(a+u)\dd u.
\]
Expanding \(f(a+u)=f(a)+u f'(a)+\cdots\), the odd term integrates to zero, while higher terms vanish with powers of \(\epsilon\).  The limit is \(f(a)\).

The Fourier representation follows from the inverse transform of the constant function \(1\):
\[
  \delta(x-a)
  =
  \int_{-\infty}^{\infty}\frac{\dd k}{2\pi}e^{\ii k(x-a)}.
\]
This equation must be read under an integral sign.  It is one of the most common formulas in QFT.

\section{Green functions}

A Green function is an inverse for a differential operator.  For
\[
  -\frac{\dd^2}{\dd x^2}G(x)=\delta(x),
\]
Fourier transform both sides.  Since \(-\dd^2/\dd x^2\) becomes \(k^2\),
\[
  k^2\tilde G(k)=1,\qquad \tilde G(k)=\frac1{k^2}.
\]
The behavior at \(k=0\) requires a boundary condition.  In three dimensions the Coulomb Green function satisfies
\[
  -\nabla^2G(\bm x)=\delta^{(3)}(\bm x),
  \qquad
  G(\bm x)=\frac{1}{4\pi|\bm x|}.
\]
Thus the solution of \(-\nabla^2\Phi=\rho\) is
\[
  \Phi(\bm x)=\int\dd^3y\,\frac{\rho(\bm y)}{4\pi|\bm x-\bm y|}.
\]
Later, the Feynman propagator will be a Green function with a quantum boundary condition.

\section*{Exercises}
\addcontentsline{toc}{section}{Exercises}

\begin{exercise}
Show that
\[
  \int_{-\infty}^{\infty}\delta(ax)f(x)\dd x=\frac1{|a|}f(0)
\]
for nonzero real \(a\).
\begin{answer}
Set \(u=ax\).  Then \(\dd x=\dd u/a\) if \(a>0\), while the limits reverse if \(a<0\); both cases give \(1/|a|\).
\end{answer}
\end{exercise}

\begin{exercise}
Using the Fourier transform convention above, prove that the transform of \(f'(x)\) is \(\ii k\tilde f(k)\), assuming boundary terms vanish.
\begin{answer}
Integrate \(\int f'(x)e^{-\ii kx}\dd x\) by parts.
\end{answer}
\end{exercise}
